\section*{Capítulo \ref{ch:g12:mechanical-waves} --- Ondas mecânicas}

\begin{solution}{exo:g12:mechanical-waves:1}
(a) $v = 320/40 = \qty{8.0}{m/s}$. (b) Espectador algum se desloca ao
longo do anel: o que viaja é a perturbação (o ato de se levantar), e não
a torcida. (c) $12/1.5 = \qty{8.0}{m/s}$ --- coerente.
\end{solution}

\begin{solution}{exo:g12:mechanical-waves:2}
(a) \omterm{def:g12:mechanical-waves:transverse}{transversal}; (b) \omterm{def:g12:mechanical-waves:transverse}{longitudinal}; (c) \omterm{def:g12:mechanical-waves:transverse}{longitudinal}; (d) \omterm{def:g12:mechanical-waves:transverse}{transversal};
(e) P \omterm{def:g12:mechanical-waves:transverse}{longitudinal}, S \omterm{def:g12:mechanical-waves:transverse}{transversal}.
\end{solution}

\begin{solution}{exo:g12:mechanical-waves:3}
(a) $d = 340 \times 4.0 \approx \qty{1.4}{km}$.
(b) $1500 \times 4.0 = \qty{6.0}{km}$. (c) A \omterm{def:g12:mechanical-waves:speed}{velocidade de propagação} no
meio atravessado.
\end{solution}

\begin{solution}{exo:g12:mechanical-waves:4}
$\lambda = v/f$: $340/440 = \qty{0.77}{m}$ no ar e
$1500/440 = \qty{3.4}{m}$ na água. A \omterm{def:g10:signals-and-waves:frequency}{frequência} fica fixa (a fonte a
dita); $v$ e $\lambda$ mudam os dois.
\end{solution}

\begin{solution}{exo:g12:mechanical-waves:5}
$\lambda = \qty{12}{cm}$; $T = \lambda/v = 0.12/0.30 = \qty{0.40}{s}$;
$f = 1/T = \qty{2.5}{Hz}$.
\end{solution}

\begin{solution}{exo:g12:mechanical-waves:6}
(a) $v = 2.4/0.30 = \qty{8.0}{m/s}$. (b) A partir da primeira foto, a
crista tem $\qty{6.0}{m}$ a percorrer: $t = 6.0/8.0 = \qty{0.75}{s}$
depois dela. (c) A fita sobe e desce, puramente na \omterm{def:g12:mechanical-waves:transverse}{transversal}, e
continua em $x = \qty{5.0}{m}$: a corda não viaja.
\end{solution}

\begin{solution}{exo:g12:mechanical-waves:7}
(a) Velocidade de aproximação \qty{4.0}{m/s} para uma separação de
\qty{4.0}{m}: $t = \qty{1.0}{s}$. (b) Superposição:
$3.0 + 2.0 = \qty{5.0}{cm}$. (c) $3.0 - 2.0 = \qty{1.0}{cm}$. (d) Dois
pulsos intactos emergem e seguem inalterados, como se nunca tivessem se
encontrado.
\end{solution}

\begin{solution}{exo:g12:mechanical-waves:8}
(a) O fino: menos massa por \omterm{def:g10:orders-of-magnitude:unit}{metro} à mesma tensão significa \omterm{def:g10:signals-and-waves:wave}{onda} mais
rápida. (b) Ela aumenta (tensão maior). (c) Não: a velocidade depende só
do meio, e não do tamanho do pulso.
\end{solution}

\begin{solution}{exo:g12:mechanical-waves:9}
Trilho: $1800/5900 \approx \qty{0.31}{s}$; ar:
$1800/340 \approx \qty{5.3}{s}$; intervalo $\approx \qty{5.0}{s}$.
\end{solution}

\begin{solution}{exo:g12:mechanical-waves:10}
Com $\lambda = \qty{24}{cm}$: (a) $2\lambda$, \omterm{def:g12:mechanical-waves:phase}{em fase};
(b) $\lambda/2$, oposição; (c) $\tfrac32\lambda$, oposição;
(d) $1.25\lambda$, nem uma coisa nem outra; (e) $\lambda/4$, nem uma
coisa nem outra.
\end{solution}

\begin{solution}{exo:g12:mechanical-waves:11}
$\lambda_P = 6.0 \times 2.0 = \qty{12}{km}$;
$\lambda_S = 3.5 \times 2.0 = \qty{7.0}{km}$ --- centenas de vezes o
tamanho de um edifício. Pontos muito mais próximos do que $\lambda$ estão
quase \omterm{def:g12:mechanical-waves:phase}{em fase}, de modo que o bairro inteiro se move junto.
\end{solution}

\begin{solution}{exo:g12:mechanical-waves:12}
(a) $T = \qty{0.50}{s}$, $f = \qty{2.0}{Hz}$.
(b) $\lambda = \qty{0.75}{m}$.
(c) $v = \lambda/T = 0.75/0.50 = \qty{1.5}{m/s}$.
(d) Em $t = T = \qty{0.50}{s}$. (e) $1.875/0.75 = 2.5$ \omterm{def:g12:mechanical-waves:wavelength}{comprimentos de
onda} $= (2 + \tfrac12)\lambda$: \omterm{def:g12:mechanical-waves:phase}{oposição de fase}.
\end{solution}

\begin{solution}{exo:g12:mechanical-waves:13}
(a) $T = 80/10 = \qty{8.0}{s}$, $f = \qty{0.125}{Hz}$.
(b) $v = \lambda/T = 120/8.0 = \qty{15}{m/s} = \qty{54}{km/h}$.
(c) $t = 900/15 = \qty{60}{s}$. (d) Não: as \omterm{def:g10:signals-and-waves:wave}{ondas} não transportam
matéria --- a boia ancorada só balança.
\end{solution}

\begin{solution}{exo:g12:mechanical-waves:14}
(a) Aproximando-se a \qty{6.0}{m/s} ao longo de \qty{6.0}{m}:
$t = \qty{1.0}{s}$; os deslocamentos opostos se cancelam --- a corda fica
momentaneamente plana. (b) Não: os pontos dela estão em movimento (as
velocidades transversais se somam, em vez de se cancelar); a \omterm{def:g10:energy-conservation:energy}{energia} é
toda cinética. (c) Os dois pulsos reaparecem e seguem inalterados. (d) A
corda plana carrega velocidade e, portanto, \omterm{def:g10:energy-conservation:energy}{energia}: a \omterm{def:g10:signals-and-waves:wave}{onda} está
momentaneamente guardada no movimento, e não na forma.
\end{solution}

\begin{solution}{exo:g12:mechanical-waves:15}
(a) $2\pi r$: \qty{3.1}{m} e \qty{50}{m} --- $\times 16$. (b) A \omterm{def:g10:energy-conservation:energy}{energia}
por \omterm{def:g10:orders-of-magnitude:unit}{metro} cai por $16$, de modo que a \omterm{def:g10:signals-and-waves:amplitude}{amplitude} diminui (a \omterm{def:g10:energy-conservation:energy}{energia} cresce
com a \omterm{def:g10:signals-and-waves:amplitude}{amplitude}). (c) As frentes retas não se alongam: sobra apenas a
absorção. (d) O meio é ativo: cada espectador se levanta com a sua
própria \omterm{def:g10:energy-conservation:energy}{energia} muscular, reabastecendo a \omterm{def:g10:signals-and-waves:wave}{onda}.
\end{solution}

\begin{solution}{pb:g12:mechanical-waves:1}
\textbf{1.} As P comprimem na direção do percurso: longitudinais; as S
sacodem de lado: transversais.

\textbf{2.} A P, mais rápida, chega primeiro em toda parte.
$t_P = 120/6.0 = \qty{20}{s}$; $t_S = 120/3.5 \approx \qty{34}{s}$.

\textbf{3.} $\Delta t = d/v_S - d/v_P = d\,(1/v_S - 1/v_P)$; para
\qty{120}{km}: $34.3 - 20 \approx \qty{14}{s}$.

\textbf{4.} Isolando $d$: $d = \dfrac{v_P v_S}{v_P - v_S}\,\Delta t
= \dfrac{6.0 \times 3.5}{2.5}\,\Delta t = \qty{8.4}{km}$ por segundo de
intervalo.

\textbf{5.} $\Delta t \propto d$: o registro de uma única estação dá a
distância dela até o epicentro sem que se conheça o instante de origem.

\textbf{6.} $d_A = 8.4 \times 25 = \qty{2.1e2}{km}$.

\textbf{7.} O intervalo dá uma distância, não uma direção: o epicentro
fica em algum ponto do círculo de raio $d_A$ em torno de A.

\textbf{8.} $d_B = 8.4 \times 15 = \qty{1.3e2}{km}$;
$d_C = 8.4 \times 32 \approx \qty{2.7e2}{km}$.

\textbf{9.} Dois círculos se encontram em dois pontos; o terceiro
seleciona um deles: a interseção comum é o epicentro.

\textbf{10.} $t_P = 210/6.0 = \qty{35}{s}$: o terremoto começou às
$09{:}14{:}32 - \qty{35}{s} = 09{:}13{:}57$.

\textbf{11.} $\qty{200}{m/s} = \qty{720}{km/h}$ --- um avião de passageiros
em velocidade de cruzeiro.

\textbf{12.} $t = \num{6.0e5}/200 = \qty{3.0e3}{s} = \qty{50}{min}$.

\textbf{13.} $t_P = 600/6.0 = \qty{100}{s}$. Aviso:
$3000 - 100 = \qty{2.9e3}{s} \approx \qty{48}{min}$.

\textbf{14.} $\lambda = vT = 200 \times 900 = \qty{1.8e5}{m} =
\qty{180}{km}$: um \omterm{def:g10:orders-of-magnitude:unit}{metro}, mais ou menos, de elevação espalhado por
\qty{180}{km} --- o navio sobe uma rampa imperceptível ao longo de vários
minutos.

\textbf{15.} $\lambda = 10 \times 900 = \qty{9.0}{km}$: a \omterm{def:g10:signals-and-waves:wave}{onda} se
comprime vinte vezes, a sua \omterm{def:g10:energy-conservation:energy}{energia} se espreme num comprimento menor, e
por isso a altura dela tem de crescer --- o tsunami se empina na costa.

\textbf{16.} $2\pi r$: \qty{63}{km} e depois $\approx \qty{1.3e3}{km}$
--- $\times 21$.

\textbf{17.} A mesma \omterm{def:g10:energy-conservation:energy}{energia} se reparte por uma frente $21$ vezes mais
longa, de modo que a \omterm{def:g10:signals-and-waves:amplitude}{amplitude} cai mesmo numa crosta sem perdas; uma
crosta real acrescenta a absorção.

\textbf{18.} As \omterm{def:g10:signals-and-waves:wave}{ondas} S, de cisalhamento \omterm{def:g12:mechanical-waves:transverse}{transversal}, não atravessam um
líquido: a sombra das \omterm{def:g10:signals-and-waves:wave}{ondas} S mostra que o \omterm{def:g11:fundamental-interactions:ladder}{núcleo} (externo) é líquido.

\textbf{19.} $\num{1e15}/\num{1.5e9} \approx \num{6.7e5}$: cerca de
setecentos mil trens a plena velocidade.

\textbf{20.} O terremoto ocorreu às $09{:}13{:}57$, no mar, no cruzamento
dos três círculos, a \qty{210}{km} da estação A. O alarme da \omterm{def:g10:signals-and-waves:wave}{onda} P do
porto tocou \qty{100}{s} depois da origem, cerca de \qty{48}{min} antes
da chegada do tsunami.
\end{solution}
